\phantomsection\label{awec:00-editors}
\marginnote{
	\begin{nscenter}
		\begin{minipage}[t]{5cm}
			\begin{center}
				\includegraphics[height=3cm]{figures/AWEC00_Editor.pdf} \\[.1cm] \textbf{Firstname1 Lastname1} \\[.05cm] University1\\[.3cm]
				\includegraphics[height=3cm]{figures/AWEC00_Editor.pdf} \\[.1cm] \textbf{Firstname2 Lastname2} \\[.05cm] University2
			\end{center}
		\end{minipage}
	\end{nscenter}
}

\begin{nscenter}
% ------------------------------------------
% TITLE SECTION
% include your title here:
{
{\Large\bfseries
% the title of your abstract goes here:
Welcome and Introduction to the Airborne Wind Energy Conference 20xx
%
}\\[2mm]
% end of the title section
%
%
%
% ------------------------------------------
% AUTHOR SECTION
% include your author(s) here:
{\bfseries
\index{Lastname1, Firstname1|textbf}
\index{Lastname2, Firstname2|textbf}
Firstname1 Lastname1$^1$, Firstname2 Lastname2$^2$
}\\
% end of the author section
%
%
%
% ------------------------------------------
% INSTITUTE SECTION
% include your institution(s) here:
$^1$University1\\
$^2$University2
}
% end of the institute section
%
\end{nscenter}


\abstractlinespacing

\begin{multicols}{2}
Dear conference participants,

Welcome to XX and welcome to the $n$th international airborne wind energy conference AWEC 20xx! We are excited to present to you an inspiring program in a beautiful location for the two conference days.

The scientific program of AWEC 20xx includes:

\begin{itemize}
\item An invited keynote presentation of 40 minutes, by

\begin{itemize}
\item  Firstname Lastname, Professor at University
\end{itemize}

\item An introductory presentation of 20 minutes, by

\begin{itemize}
\item Firstname Lastname, Professor at University
\end{itemize}

\item Three plenary presentations of 20 minutes, by

\begin{itemize}
\item Firstname Lastname, Professor at University

\item Firstname Lastname, Professor at University

\item  Firstname Lastname, Professor at University
\end{itemize}

\item $n$ contributed talk sessions in three parallel tracks with altogether $m$ presentations

\item $n$ poster sessions, each preceded by plenary spotlight presentations, with altogether $m$ poster presentations

\item $n$ panel discussions covering all aspects of airborne wind energy which include a further $m$ presentations

\end{itemize}

All abstracts presented in this book have undergone a peer review process, and we want to thank all authors and all reviewers at this place for having contributed to a high quality scientific program, as we believe.

In order to make orientation easier, we decided to rename the four main conference auditoria after renowned researchers in airborne wind energy:
\begin{itemize}
\item ``Auditorium A'' honoring Firstname Lastname (1900 -- 1970);

\item ``Auditorium B'' honoring Firstname Lastname (1910 -- 1980);

\item ``Auditorium C'' honoring Firstname Lastname (1920 -- 1990);

\end{itemize}

%%% Add photo strip

\end{multicols}

\clearpage

\phantomsection\label{photo:00-image}
\begin{tikzpicture}[remember picture, overlay, anchor=north west,inner sep=0pt, outer sep=0pt]
   \draw ([shift={(\bleed-3mm,-\bleed+3mm)}] current page.north west) node {\includegraphics[width=216mm]{figures/AWEC00_Image.pdf}};
   \captionbaselines
\end{tikzpicture}

\renewcommand*{\raggedleftmarginnote}{}
\renewcommand*{\raggedrightmarginnote}{}
\vspace*{47mm}
\marginnote{\itshape
	Image caption.
}
\vspace{-7.4mm}

\begin{multicols}{2}

%%% End photo strip

The side program of AWEC 2019 includes:
\begin{itemize}
\item a welcome reception on October 14 in the Glasgow City Chambers;

\item two lunches and four coffee breaks in the conference premises, free for all conference participants;

\item a dinner aboard the Tall Ship `Glenlee' on October 15.
\end{itemize}

The city of Glasgow is named by National Geographic as one of its ``Best of the World'' destinations, while voted by Rough Guide readers the world’s friendliest city! Glasgow is a city with a very strong, indeed a globally renowned, knowledge base sector and vibrant wind energy sector. UK’s largest onshore wind farm, Whitelee, is just 20 minutes from the city centre.

Founded in 1796 as the Andersonian Institute to be a ``place of useful learning'', the University of Strathclyde received its royal charter in 1964 as the UK’s first technological university. Based right in the very heart of Glasgow, the University of Strathclyde was awarded Scottish University of the Year 2020 by the Times and Sunday Times Good University Guide.

The Wind Energy and Control Centre (WECC), in the Department of Electronic \& Electrical Engineering at the University of Strathclyde, is one of the largest wind energy research groups in the world with over 80 research assistants and PhD students. WECC has expertise in turbine and powertrain design; fault diagnosis, failure rate analysis, O\&M and asset management; offshore networks, connection-to-shore and grid integration; power production forecasting, turbine and array dynamics, modelling and simulation; turbine and wind farm control.

WECC leads the UK’s pre-eminent doctoral training programme in offshore renewable energy, the EPSRC Centre for Doctoral Training in Wind \& Marine Energy Systems \& Structures (CDT-WAMSS). The centre was first established in 2009 and over the past decade it has trained more than 100 doctoral students, working with over 40 industrial partners and with graduates providing expertise across all areas of wind and marine energy engineering. As of October 2019, the CDT brings together the leading UK research groups in Wind Energy at Strathclyde, Marine Energy at the University of Edinburgh and Offshore Structures at the University of Oxford.

\end{multicols}
\clearpage

\begin{multicols}{2}

%Since 2007, the University has a Center for Renewable Energies (ZEE), which organizes in particular an international master program in Renewable Energy Engineering and Management (REM). Finally, in 2016, the university has founded a new Department of Sustainable Systems Engineering (INATECH) within the Faculty of Engineering and fostered an alliance with the Fraunhofer society within the new ``Sustainability Center Freiburg'', which cosponsors AWEC 2017. Since 2015, Freiburg University is also part of the European Initial Training Network AWESCO on Airborne Wind Energy (coordinated by TU Delft) which also cosponsors AWEC 2017.

The conference would not have been possible without the support of its sponsors, who are listed on pages \pageref{awec:sponsors}--\pageref{awec:sponsorsend} and, to which we want to express our sincere gratitude. We are also grateful to the City of Glasgow and the Lord Provost for hosting the Welcome Reception in the Glasgow City Chambers. A special thanks also goes to the Glasgow Convention Bureau and the TIC conference staff for providing exceptional support to this event.

We also want to thank all members of the programme committee and organising committee -- listed on page \pageref{awec:committees} -- for their efforts in making the conference a success. And within the organising committee, we want in particular to thank, Stefanie Thoms for her outstanding contributions that have not only made this conference possible but have greatly enhanced the conference experience for all participants.

Last but not least, we are grateful to you, the participants of AWEC 2019, not only for coming to the conference, but also for your various contributions and your hopefully active participation in the discussions during pannel sessions, after talks, at lunches, dinners and coffee breaks. We very much look forward to an inspiring and exciting conference together with you!

Sincerely,
\end{multicols}

\regularlinespacing

\vspace{1mm}

\begin{minipage}[t]{.3\textwidth}
	\raisebox{5mm}{\includegraphics[width=3.5cm]{figures/signature1.pdf}}
\end{minipage}
\begin{minipage}[t]{.3\textwidth}
    \raisebox{5mm}{\includegraphics[width=3.5cm]{figures/signature2.pdf}}
\end{minipage}

\vspace{1mm}

\begin{minipage}[t]{.3\textwidth}
    Firstname1 Lastname1\\[.1cm]
    University1\\
    City1, Country1
\end{minipage}
\begin{minipage}[t]{.3\textwidth}
    Firstname2 Lastname2\\[.1cm]
    University2 \\
    City2, Country2
\end{minipage}

\newpage
